\section{Resultados Esperados}

De acuerdo con las brechas identificadas en el diagnóstico inicial de madurez digital y los objetivos estratégicos definidos para el proyecto de Industria Textil Atlas E.I.R.L., la implementación del sistema Odoo se orientó a lograr los siguientes resultados esperados en la cadena logística y comercial:

\begin{itemize}
    \item \textbf{Ordenamiento Logístico y Control de Almacén:} Se proyectó la centralización definitiva del catálogo de 22 artículos de protección y accesorios deportivos utilizando la codificación SKU de referencia interna. Esto permite eliminar las discrepancias físicas y el desorden, asegurando que las existencias reales coincidan con los datos de almacén.
    \item \textbf{Reducción del Tiempo de Atención en Tienda:} El resultado esperado más directo es la agilización del despacho en el mostrador de San Camilo. Mediante la terminal del POS con interfaz de búsqueda rápida por categorías, se planificó reducir el tiempo promedio de atención y preparación de pedidos de 15--20 minutos a menos de 5 minutos por cliente.
    \item \textbf{Aprovisionamiento Automatizado y Reglas de Reorden:} Se proyectó transitar de un modelo de compras reactivas hacia un sistema de reposición inteligente. Al configurar reglas de stock mínimo automático para la mercadería tercerizada, se esperaba que el sistema sugiriera de manera proactiva solicitudes de compra a los talleres externos cuando el stock descienda de los niveles de seguridad.
    \item \textbf{Control de Costos y Márgenes Gerenciales:} Se esperaba dotar al gerente general Víctor Huayllaro de una interfaz de control financiero exclusiva. A través del acceso gerencial a la ficha de producto y al panel de valoración de inventarios, se buscó transparentar los costos reales de adquisición de talleres y simular márgenes comerciales exactos en soles (PEN) para fijar precios competitivos sin recurrir al cálculo intuitivo o manual.
    \item \textbf{Agilidad Comercial Mayorista y Presupuestos Rápidos:} A través del módulo de Ventas, se proyectó digitalizar la generación de propuestas comerciales para colegios y clubes deportivos. El personal de ventas debía ser capaz de generar presupuestos profesionales en formato PDF y compartirlos de manera inmediata por WhatsApp, eliminando la duplicación de tareas administrativas.
\end{itemize}
